\section{Feasible Inference}
\label{sec:clt}

\subsection{Oracle Weighted Least Squares}

With known weights $w_j^* = 1/v_j = n\,D_j^2/\sigma_0^2$, the oracle
WLS estimator of $\boldsymbol{\beta} = (\alpha, H)^\top$ in
\eqref{eq:log-model} is
\begin{equation}
    \boldsymbol{\hat\beta}^* = (\mathbf{X}^\top \mathbf{W}^* \mathbf{X})^{-1}
    \mathbf{X}^\top \mathbf{W}^* \mathbf{Y},
    \label{eq:oracle-wls}
\end{equation}
where $\mathbf{X} = [\mathbf{1}\mid \mathbf{x}]$, $\mathbf{x} = (\log
1, \ldots, \log L)^\top$, and $\mathbf{W}^* = \diag(w_1^*, \ldots, w_L^*)$.

\begin{proposition}[Asymptotic Normality of the Oracle WLS Estimator]
\label{prop:oracle-clt}
Under \Cref{ass:signal,ass:sigma,ass:lags}, as $n \to \infty$ with $L$
fixed,
\[
    n^{1/2}\,\bigl(\hat H^* - H\bigr) \;\dto\; \calN\!\left(0,\;
    \sigma_0^2 \bigl[(\mathbf{X}^\top \mathbf{W}^* \mathbf{X})^{-1}\bigr]_{22}\right),
\]
where $[\,\cdot\,]_{22}$ denotes the $(2,2)$ entry.
Explicitly,
\begin{equation}
    \Var(\hat H^*) \;\approx\; \frac{\sigma_0^2}{n}
    \cdot \frac{S_{w^*}}
    {S_{w^*} S_{w^* x^2} - (S_{w^* x})^2},
    \label{eq:oracle-var}
\end{equation}
where $S_{w^*} = \sum_j w_j^*$, $S_{w^* x} = \sum_j w_j^* x_j$,
$S_{w^* x^2} = \sum_j w_j^* x_j^2$.
\end{proposition}

\begin{proof}
Standard weighted least squares with $\mathbf{U} \sim
\calN(\mathbf{0}, \sigma_0^2 n^{-1} (\mathbf{W}^*)^{-1})$; the
$(2,2)$ entry of $(\mathbf{X}^\top \mathbf{W}^* \mathbf{X})^{-1}$ follows
by $2\times 2$ matrix inversion.
\end{proof}

\begin{remark}[Lag-energy information law]
    For the power-law family, the oracle weights satisfy
    $w_j^* = (n\zeta^2/\sigma_0^2) j^{2H}$.  Hence the scalar lag-energy
    $\mathcal I_{n,L}(H) = n\sum_{j=1}^L j^{2H}$ is the basic information
    budget behind the matrix $\mathbf X^\top \mathbf W^* \mathbf X$.
    This same quantity reappears in the lower bound and in the departure scale
    of the adequacy statistic under local alternatives.
\end{remark}

\subsection{Two-step feasible WLS}

Since $D_j$ is unknown, the oracle weights are infeasible. We use the
following two-step procedure.

\begin{algorithm}[H]
\caption{Two-Step FWLS Estimator}\label{alg:fwls}
\begin{algorithmic}[1]
\State \textbf{Input:} $(\Dhj)_{j=1}^L$, $n$, and $\sigma_0$.
\State \textbf{Step 1 (Pilot OLS).} Compute $(\hat\alpha_{\mathrm{pil}}, \hat H_{\mathrm{pil}})$ by unweighted OLS of $\log \Dhj$ on $(\mathbf{1}, \log j)$.
\State \textbf{Step 2 (Imputation).} Set $\hat D_j^{\mathrm{pil}} = \exp(\hat\alpha_{\mathrm{pil}} + \hat H_{\mathrm{pil}} \log j)$.
\State \textbf{Step 3 (FWLS).} Set $\hat w_j = n (\hat D_j^{\mathrm{pil}})^2 / \sigma_0^2$ and compute $(\hat\alpha, \Hhat)$ by WLS with weights $(\hat w_j)$.
\State \textbf{Output:} $\Hhat$, $\hat\alpha$, and residuals $\hat U_j = \log \Dhj - \hat\alpha - \Hhat \log j$.
\end{algorithmic}
\end{algorithm}

\begin{theorem}[Asymptotic Normality of the Feasible WLS Estimator]
\label{thm:clt-fwls}
Under \Cref{ass:signal,ass:sigma,ass:lags}, the two-step FWLS estimator
$\Hhat$ satisfies, as $n \to \infty$ with $L$ fixed:
\begin{equation}
    n^{1/2}\bigl(\Hhat - H\bigr) \;\dto\;
    \calN\!\Bigl(0,\; \sigma_0^2
    \bigl[(\mathbf{X}^\top \mathbf{W}^* \mathbf{X})^{-1}\bigr]_{22}\Bigr).
    \label{eq:fwls-clt}
\end{equation}
In particular, $\Hhat$ is \emph{asymptotically equivalent to the oracle
WLS estimator}: the pilot estimation error inflates neither the asymptotic
variance nor the rate.
\end{theorem}

\begin{corollary}[Joint law and smooth downstream functionals]
\label{cor:joint-law-smooth-functional}
Under \Cref{ass:signal,ass:sigma,ass:lags},
\[
  n^{1/2}
  \begin{pmatrix}
    \hat\alpha-\alpha \\
    \Hhat-H
  \end{pmatrix}
  \dto
  \calN\!\left(
    \mathbf 0,
    \frac{\sigma_0^2}{\zeta^2\Delta_H}
    \begin{pmatrix}
      R_2(H) & -R_1(H) \\
      -R_1(H) & R_0(H)
    \end{pmatrix}
  \right),
\]
where
\[
  R_k(H) = \sum_{j=1}^L j^{2H}(\log j)^k,
  \qquad
  \Delta_H = R_0(H)R_2(H)-R_1(H)^2.
\]
Equivalently,
\[
  \sqrt{\mathcal I_{n,L}(H)}
  \begin{pmatrix}
    \hat\alpha-\alpha \\
    \Hhat-H
  \end{pmatrix}
  \dto \calN(\mathbf 0,\Sigma_{L,H,\zeta}),
\]
with
\[
  \Sigma_{L,H,\zeta}
  =
  \frac{\sigma_0^2R_0(H)}{\zeta^2\Delta_H}
  \begin{pmatrix}
    R_2(H) & -R_1(H) \\
    -R_1(H) & R_0(H)
  \end{pmatrix}.
\]
Therefore, for any $C^1$ map $g$ defined on a neighborhood of
$(\alpha,H)$,
\[
  \sqrt{\mathcal I_{n,L}(H)}
  \bigl(g(\hat\alpha,\Hhat)-g(\alpha,H)\bigr)
  \dto
  \calN\!\left(
    0,
    \nabla g(\alpha,H)^\top \Sigma_{L,H,\zeta}\nabla g(\alpha,H)
  \right).
\]
\end{corollary}

\begin{proof}
The weighted least-squares normal equations are two-dimensional. Inverting the
oracle information matrix gives the displayed covariance formula, and the same
perturbation argument as in \Cref{thm:clt-fwls} transfers the law from oracle WLS
to FWLS. The final claim is the multivariate delta method.
\end{proof}

\begin{remark}[$\sigma_0$ unknown]
    When $\sigma_0$ is unknown, a pilot plug-in estimate constructed from the
    residual scale yields a self-normalized FWLS fit. In the adequacy test,
    a naive $\chi^2(L-3)$ calibration is strongly conservative because the
    estimated scale and the residual statistic are coupled through the same
    data. Under sample splitting, the Gaussian benchmark with oracle profile and
    independent scale estimate has the exact split-$F$ law of
    \Cref{prop:oracle-split-f}, whereas the fully feasible split procedure
    remains conservative on the tested grid. For the same-sample procedure, a
    parametric bootstrap based on the fitted null profile restores near-nominal
    size; see \Cref{tab:sigma0-plugin}.
\end{remark}

\begin{proof}[Proof sketch]
Write $\hat w_j = w_j^* \cdot r_j$ where $r_j = (\hat D_j^{\mathrm{pil}}
/ D_j)^2$.  Since $\hat H_{\mathrm{pil}} - H = O_p(n^{-1/2})$ (OLS
rate in $\log$-scale with bounded regressors), we have
$\log r_j = 2(\hat H_{\mathrm{pil}} - H)\log j + O_p(n^{-1}) = O_p(n^{-1/2}\log j)$,
hence $r_j = 1 + O_p(n^{-1/2}\log j)$ uniformly in $j$.

The FWLS estimator with weights $\hat w_j = w_j^*(1 + \delta_j)$,
$\delta_j = r_j - 1 = O_p(n^{-1/2}\log j)$, satisfies
\[
    \Hhat - H = \hat H^* + R_n,
\]
where the remainder $R_n = O_p(n^{-1}\log L)$ via a second-order
Taylor expansion of the WLS normal equations in the weight perturbation.
Since $n^{1/2} R_n = O_p(n^{-1/2}\log L) \to 0$, the Slutsky theorem
gives \eqref{eq:fwls-clt}.
\end{proof}
